\section[Demonstrations]{\secstyle{Demonstrations:}}

\subsection[Unicité de la limite]{\subsecstyle{Si \((u_n)\) converge vers \(l\) et \(l'\), alors \(l = l'\):}}

\begin{proof}[\unskip\nopunct]
   Soit \((u_n)\) une suite qui converge vers \(l\) et \(l' \in \R\), supposons que \(l \neq l'\) et donc que \(|l - l'| \neq 0\).\<
   
   Soit \(\epsilon\) strictement positif, alors on a:
   \[
      \begin{cases}
         \exists N_1 \in \N \; , \; \forall n > N_1 \; ; \; |u_n - l| < \epsilon \\ 
         \exists N_2 \in \N \; , \; \forall n > N_2 \; ; \; |u_n - l'| < \epsilon
      \end{cases} 
   \]
   
   Posons, \(M := \max(N_1, N_2)\), soit \(n > M\). On remarque que les deux propriétés ci-dessus sont vraies à partir du rang \(M\).
   Sommons les deux inégalités précédentes pour obtenir:
   \[
      |u_n - l| + |u_n - l'|< 2\epsilon
   \]
   Les propriétés des inégalités nous donnent:
   \[
      |l - l'| \leq |l - u_n| + |u_n - l'|< 2\epsilon   
   \]
   Et donc, on obtient:
   \[
      |l - l'| < 2\epsilon
   \]
   Or, en posant \(\epsilon := \frac{|l - l'|}{2}\), l'inégalité est absurde, et donc nécessairement, on a \(|l - l'| = 0\) et donc \(l = l'\) 

\end{proof}

\subsection[Somme de suites convergentes]{\subsecstyle{Montrer que la suite \((u_n+v_n)\) converge vers \(l + l'\):}}
\begin{proof}[\unskip\nopunct]
   Soient \((u_n)\) et \((v_n)\) deux suites qui convergent vers \(l, l' \in \R\) respectivement, soit \(\epsilon\) strictement positif, on sait que:
   \[
   \begin{cases}
         \exists N_1 \; , \; \forall n > N_1 \; ; \; |u_n - l| < \epsilon\\
         \exists N_2 \; , \; \forall n > N_2 \; ; \; |v_n - l'| < \epsilon\\
   \end{cases}
   \]
   Posons \(M := \max(N_1, N_2)\), soit \(n > M\), on a:
   \[
   \begin{cases}
         M \geq N_1 \text{ ET } |u_n - l| < \epsilon\\
         M \geq N_2 \text{ ET }|v_n - l'| < \epsilon\\
   \end{cases}
   \]
   Puis par somme et utilisation de l'inégalité triangulaire:
   \[
      |(u_n + v_n) - (l + l')| \leq |u_n - l| + |v_n - l'| < 2 \epsilon
   \]
   Quitte à prendre un \(\epsilon\) deux fois plus petit, on a donc:
   \[
      \forall \epsilon > 0 \; , \; \forall n > M \; ; \; |(u_n + v_n) - (l + l')| < \epsilon
   \]
   On en conclut que la somme des suites converge bien vers la somme des limites.
   
\end{proof}

\subsection[Caractérisation des suites croissantes]{\subsecstyle{\textbf{Caractérisation des suites croissantes:}}}
Démontrons \(\P\) l'équivalence:
\[
\P := (\forall n,m \in \N\; ; \; n \leq m \implies u_n \leq u_m) \Longleftrightarrow (\forall n \in \N\; ; \; u_n \leq u_{n+1})
\]
Posons \(\P_1\) le membre de gauche et \(\P_2\) le membre de droite de \(\P\) pour faciliter la lecture.

\begin{proof}[\unskip\nopunct]
\boxed{\implies} Supposons \(\P_1\): \(\forall n,m \in \N\; ; \; n \leq m \implies u_n \leq u_m\)\<

Soit \(n \in \N\) quelconque. On sait que \(\P_1\) est vraie pour tout \(m\) donc en particulier, elle est vraie pour \(m := n + 1 \in \N\), qui est bien supérieur à \(n\). Posons \(m := n + 1\).

Donc on a:
\[
   n \leq m \implies u_n \leq u_{n + 1}
\]
La propriété étant valable pour tout \(n \in \N\), on en déduit \(\P_2\).\<

\boxed{\impliedby} Supposons \(\P_2\): \(\forall n \in \N\; ; \; u_n \leq u_{n+1}\)\< 

Soit \(p, q \in \N\) tels que \(p \leq q\), montrons que \(u_p \leq u_q\), alors d'aprés nos hypothèses, on a:

\[
   (u_p \leq u_{p+1}) \text{ et } (u_{p+1} \leq u_{p+2})\ldots\ldots \text{ et } (u_{q-2} \leq u_{q-1}) \text{ et } (u_{q-1} \leq u_{q})    
\]
On en déduit \(\P_1\) par transitivité de la relation \(\leq\).
\end{proof}

\subsection[Toute suite bornée est majorée et minorée]{\subsecstyle{\textbf{Toute suite est bornée ssi elle est majorée et minorée:}}}
Soit \((u_n)\) une suite quelconque et \(n \in \N\), prouvons l'équivalence suivante notée \(\P\):
\[
   \P:= (\exists M \in \R \; , \; \forall n \in \N \; ; \; |u_n| \leq M) \Longleftrightarrow \Bigr[(\exists M \in \R \; , \; \forall n \in \N \; ; \; u_n \leq M) \text{ ET } (\exists m \in \R \; , \; \forall n \in \N \; ; \; u_n \geq m)\Bigr]
\]

Démontrons \(\P\) par double implication:

\begin{proof}[\unskip\nopunct]
   \boxed{\implies} Supposons que \((u_n)\) est bornée, alors il existe un \(K \in \R\) tel qu'on ait:
   \[
      |u_n| \leq K
   \]
   Par définition de la valeur absolue, on a l'équivalence:
   \[
      -K \leq u_n \leq K
   \]
   Et donc on a bien trouvé un minorant et un majorant de \((u_n)\) dans \(\R\) qui sont respectivement:
   \begin{alignat*}{2}
      m &:= -K \\
      M &:= K
   \end{alignat*}

   \boxed{\implies} Supposons que \((u_n)\) soit majorée et minorée et posons \(m\) son minorant et \(M\) son majorant, alors par définition, on a:
   \[
      m \leq u_n \leq M
   \]
   Montrons que \((u_n)\) est bornée, posons \(\lambda := \max(|m|, |M|)\), alors on a successivement:
   \[
      -\lambda \leq - |m| \leq \boxed{m \leq u_n \leq M} \leq |M| \leq \lambda
   \]
   Puis par transitivité:
   \[
      -\lambda \leq u_n \leq \lambda      
   \]
   Et donc par définition de la valeur absolue, on a:
   \[
      |u_n| \leq \lambda
   \]
   On a donc bien trouvé un \(\lambda \in \R\) tel que \((u_n)\) soit bornée par ce \(\lambda\).
\end{proof}

\subsection[Toute suite périodique est bornée]{\subsecstyle{\textbf{Toute suite périodique est bornée:}}}

Soit \((u_n)\) une suite périodique quelconque et \(n \in \N\). Montrons que \((u_n)\) est bornée, plus précisément:

\[
   (\exists M \in \R \;,\; \forall n \in \N \; ; \; (u_n) \leq M) \text{ et } (\exists m \in \R \;,\; \forall n \in \N \; ; \; (u_n) \geq m)
\]

\begin{proof}[\unskip\nopunct]
   On sait que:
   \[
      \exists T \in \N^* \; , \; \forall N \in \N \; ; \; u_{n+T} = u_n
   \]
   Donc on a:
   \[
      u_n \in E:=\Bigl\{ u_0,u_1, \ldots,u_{T-2}, u_{T-1}\Bigl\}
   \]
   Or, E est un ensemble fini et tout ensemble fini possède un minimum est un maximum.

   Donc en posant \(M := \max(E)\) et \(m := \min(E)\), on a:
   \[
      m \leq u_n \leq M   
   \]
   \((u_n)\) est donc majorée et minorée, et on en conclut qu'elle est bornée.
\end{proof}
\pagebreak

\section[Exercices]{\secstyle{Exercices:}}
